% Les annexes du guide.
%
% Le code qu'elles présentent n'est pas saisi ici : docs/gendoc.py
% l'extrait des paquets à chaque composition, dans docs/generated/.
% Ce fichier ne contient que la prose et l'ordre de présentation.
%
% Ajouter un exemple demande donc deux gestes : la fonction Example
% dans example_test.go, et le \input correspondant ci-dessous. Si le
% second manque, l'exemple n'est pas dans le guide ; s'il désigne un
% exemple disparu, la composition s'arrête et le dit.

\appendix

\chapter{Exemples d'utilisation}\label{chap:exemples}

Les exemples qui suivent ne sont pas des extraits recopiés. Ce sont des
fonctions \texttt{Example} des deux paquets : \texttt{go\ test} les
compile, les exécute, et compare ce qu'elles écrivent au commentaire
\texttt{//\ Output} qui les termine. Un exemple qui ne correspond plus
au code fait échouer les tests.

C'est la seule forme de documentation qui ne puisse pas vieillir sans
qu'on le sache, et la raison pour laquelle ce qui suit est extrait de
\texttt{example\_test.go} au moment de composer le guide plutôt que
saisi une fois pour toutes. Les mêmes exemples se déplient sur
pkg.go.dev, sous la fonction qu'ils illustrent.

Ils sont écrits dans le paquet de test externe : ils montrent donc
l'import et le préfixe qu'un appelant écrit réellement, et non les
raccourcis que seul le paquet lui-même s'autorise. Les commentaires
sont en anglais, comme le code.

Deux fonctions auxiliaires, \texttt{at} et \texttt{show}, abrègent les
exemples de \texttt{timeseries} : la première fabrique un relevé pris
\texttt{n} minutes après 10:00 UTC, la seconde affiche une série à
raison d'une ligne par relevé.

\section{\texttt{notavalue}}

\subsection{Les trois règles sur un même calcul}

Un relevé manquant est ignoré, un calcul rompu se propage, et un
ensemble où il n'y a rien à calculer ne donne pas zéro mais rien.
Chapitre~\ref{chap:nav}.

\input{generated/exemple-notavalue-Example}

\subsection{L'asymétrie de l'addition}

L'addition traite le manque comme neutre, les trois autres opérations
non. Une pluie mensuelle dont un jour n'a pas été relevé reste la pluie
tombée ; une concentration manquante multipliée par un débit connu ne
donne rien.

\input{generated/exemple-notavalue-ExampleAdd}

\subsection{Choisir la première source disponible}

Un capteur principal, son secours, puis une valeur modélisée.
\texttt{Coalesce} passe sur ce qui manque, mais pas sur une erreur : un
calcul rompu en amont ne se répare pas en passant au plan B.

\input{generated/exemple-notavalue-ExampleCoalesce}

\subsection{Ce que vaut un résultat}

Une moyenne ne dit pas sur combien de relevés elle repose. Les
compteurs le disent, et distinguent les relevés absents des calculs
rompus --- ces derniers sont à chercher en amont.

\input{generated/exemple-notavalue-ExampleCountUsable}

\subsection{Garder la distinction à l'écran}

Les verbes de la bibliothèque standard voient les deux sentinelles
comme des NaN et écrivent « NaN » pour l'une comme pour l'autre. Toute
la distinction disparaît dans une ligne de journal ou un message
d'échec de test ; \texttt{Format} la conserve.
Chapitre~\ref{chap:sorties}.

\input{generated/exemple-notavalue-ExampleFormat}

\subsection{Une médiane qui n'invente pas de valeur}

Sur un nombre pair de relevés, cette médiane rend la plus petite des
deux valeurs centrales plutôt que leur moyenne. Sur une vanne codée 0
et 1, la définition classique renverrait 0,5 --- un état que
l'équipement n'a jamais occupé.

\input{generated/exemple-notavalue-ExampleMedian}

\section{\texttt{timeseries}}

\subsection{Le parcours complet}

Des relevés bruts à un résumé : les valeurs invraisemblables écartées,
une grille régulière, les trous courts comblés. Les relevés sont remis
dans n'importe quel ordre, ce qui ne coûte rien --- chapitre
\ref{chap:serie}.

\input{generated/exemple-timeseries-Example}

\subsection{Un trou et une erreur ne sont pas le même accident}

La première série a un relevé manquant, et garde une moyenne. La
seconde porte le résultat d'un calcul rompu, et n'en a plus. La même
asymétrie se lit sur la variation d'un point au suivant : un relevé de
20 après un trou n'est pas une hausse de 20.
Chapitre~\ref{chap:resume-stats}.

\input{generated/exemple-timeseries-Example-gapVersusError}

\subsection{Nettoyer, et garder ce qu'on a rejeté}

Le nettoyage rend deux séries : celle qui est nettoyée et celle des
relevés écartés, pour que la décision puisse être revue. Dans la série
nettoyée, l'instant rejeté est toujours là --- sa mesure est devenue
manquante, elle n'a pas disparu. Chapitre~\ref{chap:nettoyage}.

\input{generated/exemple-timeseries-ExampleTimeSeries-RemoveOutbounds}

\subsection{Un enregistreur qui dérive}

Un enregistreur censé rapporter à l'heure juste rapporte quelques
secondes trop tard. Sans tolérance, chacun de ces relevés tombe dans la
fenêtre suivante, ce qui laisse une fenêtre vide et met deux relevés
dans la suivante. Chapitre~\ref{chap:regularisation}.

\input{generated/exemple-timeseries-ExampleTimeSeries-RegularizeWithTolerance}

\subsection{Combler un silence court, pas une panne}

La limite se mesure entre les deux relevés qui encadrent le trou, et
non entre les pas de la grille : elle veut donc dire la même chose sur
une série régulière et sur une série brute.
Chapitre~\ref{chap:interpolation}.

\input{generated/exemple-timeseries-ExampleTimeSeries-InterpolateWithin}

\subsection{Comprimer un signal d'état, puis le restituer}

Un signal qui passe son temps à ne pas bouger. \texttt{Reduce} garde
les instants où quelque chose s'est produit, \texttt{Expand} remet la
grille. Chapitre~\ref{chap:reduction}.

\input{generated/exemple-timeseries-ExampleTimeSeries-Reduce}

\subsection{Parcourir un conteneur sans allouer}

Un conteneur tient les séries qui vont ensemble --- les relevés bruts,
la version nettoyée, chaque variante produite en chemin.
\texttt{MeasTo} est le compagnon de la boucle intérieure : on lui
confie une tranche de longueur nulle, il la remplit au lieu d'en
allouer une neuve, et un passage sur cent séries n'alloue qu'une fois.
Chapitre~\ref{chap:conteneurs}.

\input{generated/exemple-timeseries-ExampleTsContainer}

\chapter{Référence des interfaces}\label{chap:api}

Les listes qui suivent donnent, pour chaque paquet, ce qu'il exporte :
constantes, variables, fonctions, types avec leurs champs et leurs
méthodes. Elles sont extraites de \texttt{go\ doc\ -all} et ne peuvent
donc décrire ni une fonction qui n'existe plus, ni oublier celle qui
vient d'être ajoutée.

Ce sont des signatures, sans les commentaires qui les accompagnent. Le
godoc complet est sur pkg.go.dev, pour
\href{https://pkg.go.dev/github.com/fflamingodev/notavalue}{\texttt{notavalue}}~\cite{pkgnotavalue}
et pour
\href{https://pkg.go.dev/usefulrisk.com/timeseries}{\texttt{timeseries}}~\cite{pkgtimeseries},
où il est navigable et à jour de la version publiée. Le reproduire ici
doublerait ce document pour répéter en anglais ce que les chapitres
précédents expliquent en français.

L'ordre est celui de \texttt{go\ doc} : alphabétique à l'intérieur de
chaque rubrique, les constructeurs et les méthodes d'un type venant
après sa déclaration.

\input{generated/api-notavalue}

\input{generated/api-timeseries}
